\begin{lemma}
\label{editorial-source-lemma-image-stable-specialization-closed}
Let $R \to S$ be a ring map. Let $T \subset \Spec(R)$
be the image of $\Spec(S)$. If $T$ is stable under specialization,
then $T$ is closed. More precisely, for every ring map
$\varphi:R\to S$, without a specialization hypothesis,
$$
\overline T=V(\ker\varphi)
=\{\mathfrak p:\text{there is }\mathfrak p'\in T
\text{ with }\mathfrak p'\subseteq\mathfrak p\}.
$$
\end{lemma}

\begin{proof}
We retain both original proofs and all original rings.
Write the map as $\varphi:R\to S$ and put $I=\ker\varphi$.

\medskip\noindent
First proof. Let $\mathfrak p$ be a point of $\overline T$.
Every $f\notin\mathfrak p$ gives $D(f)\cap T\ne\varnothing$.
The prime localization correspondence identifies this
intersection with the image of $\Spec(S_{\varphi(f)})$,
so $S_{\varphi(f)}\ne0$.
The directed localization comparison with $S_{\mathfrak p}$
can be written without suppressing the original coefficient map.
Index finite subsets $E\subseteq R\setminus\mathfrak p$
by inclusion and put $f_E=\prod_{f\in E}f$.
For $E\subseteq F$, the transition is
$$
S_{\varphi(f_E)}\longrightarrow S_{\varphi(f_F)},\qquad
\frac{s}{\varphi(f_E)^n}\longmapsto
\frac{s\varphi(f_{F\setminus E})^n}{\varphi(f_F)^n}.
$$
The colimit map to $S_{\mathfrak p}=(R\setminus\mathfrak p)^{-1}S$
sends a fraction to the same fraction. It is surjective
because $s/\varphi(t)$ comes from $E=\{t\}$.
If an original numerator $s$ maps to zero, some
$u\notin\mathfrak p$ has $\varphi(u)s=0$; the numerator
then vanishes at the stage $E\cup\{u\}$ where
$\varphi(u)$ is a unit. Thus the map is also injective.
Every stage is nonzero, including $E=\varnothing$,
because $f_E\notin\mathfrak p$; their common unit
cannot die in the colimit, since that would make it
zero at a nonzero later stage. Hence $S_{\mathfrak p}\ne0$.

Choose a prime of this nonzero ring and contract it
to a prime $\mathfrak q\subset S$. All
$\varphi(t)$, $t\notin\mathfrak p$, are units in
$S_{\mathfrak p}$, so
$\mathfrak p'=\varphi^{-1}(\mathfrak q)\subseteq\mathfrak p$.
This is a point of the original image $T$.
Thus every point of $\overline T$ is a specialization
of a point of $T$. The reverse containment follows
because $\overline T$ is closed and contains $T$.

The same localization argument identifies the closed set
exactly. Every prime in $T$ contains $I$, so
$\overline T\subseteq V(I)$. Conversely, for
$I\subseteq\mathfrak p$, the ring $S_{\mathfrak p}$
cannot be zero: the equality $1/1=0$ would give
$t\notin\mathfrak p$ with $\varphi(t)1=0$,
contradicting $t\in I\subseteq\mathfrak p$.
Its prime again contracts to a point of $T$ contained
in $\mathfrak p$. This proves both displayed equalities.
If $T$ is stable under specialization, they imply
$T=V(I)$ and hence the asserted closedness.

\medskip\noindent
Second proof. Retain the original quotient map
$q:R\to R/I$ and the induced injective map
$\overline\varphi:R/I\to S$, $r+I\mapsto\varphi(r)$.
The spectrum map for $q$ is the homeomorphism onto
$V(I)$ given by contraction, with inverse
$\mathfrak p\mapsto\mathfrak p/I$.
For each $\mathfrak p\in V(I)$ choose a minimal prime
$\overline{\mathfrak p}_0\subseteq\mathfrak p/I$
of $R/I$. Existence follows from Zorn's lemma: the
intersection of a chain of primes contained in
$\mathfrak p/I$ is prime and is a lower bound.
Indeed if $ab$ belongs to that intersection and $a$
is absent from one chain member, primality forces
$b$ into that member and all smaller members, and
then into all larger members as well.
Lemma \ref{algebra-lemma-injective-minimal-primes-in-image}
gives a prime $\mathfrak q\subset S$ whose contraction
along $\overline\varphi$ is $\overline{\mathfrak p}_0$.
Along the original map its contraction is
$$
\varphi^{-1}(\mathfrak q)
=q^{-1}(\overline{\mathfrak p}_0)\subseteq\mathfrak p.
$$
Thus every point of $V(I)$ is again a specialization
of an original image point; no original ring has been
replaced in the comparison. If $S=0$, then $I=R$,
and all three sets in the statement are empty.
\end{proof}